\documentclass{beamer}

\usetheme[compactlogo,white]{Madison}

\title{Modeling Cognitive Load as an Instructional Partitioning Problem}
\subtitle{EDPS 506 Peer Review Proposal}
\author{Yifan Yang}
\date{Spring 2026}

\begin{document}

\small

\begin{frame}[t]{Background and Research Questions}

\textit{A theory-first project on cognitive load, segmentation, and learning under working-memory constraints.}

\vspace{0.4em}
\textbf{Background}
\begin{itemize}
\setlength{\itemsep}{0.15em}
\item Cognitive load theory explains why working-memory limits matter for learning.
\item Current theory is under-formalized at the level of structure, segmentation, and outcomes.
\item Segmentation can ease local burden but weaken coherence needed for transfer.
\end{itemize}

\vspace{0.2em}
\textbf{Research Questions}
\begin{itemize}
\setlength{\itemsep}{0.15em}
\item How should structured instructional information be modeled under working-memory constraints?
\item How do segmentation and representation shape the trade-off between local relief and preserved coherence?
\item Can that trade-off predict recall and transfer?
\end{itemize}

\end{frame}

\begin{frame}[t]{Proposed Model}

\begin{block}{Core Idea}
\begin{itemize}
\setlength{\itemsep}{0.15em}
\item Learning material is represented as interconnected information elements.
\item Intrinsic cost reflects information quantity, element complexity, and coordination demand.
\item Extraneous cost reflects how the material is represented or presented.
\item Segmentation can lower local burden but also increase fragmentation across parts.
\end{itemize}
\end{block}

\begin{alertblock}{Main Mechanism}
\begin{itemize}
\setlength{\itemsep}{0.15em}
\item Goal-relevant processing is what remains after intrinsic and extraneous costs are subtracted from limited working-memory capacity.
\item Learning outcomes depend on both available processing and preserved relational structure.
\item Recall should depend more on available processing; transfer should be more sensitive to fragmentation.
\end{itemize}
\end{alertblock}

\end{frame}

\begin{frame}[t]{Hypotheses and Validation Path}

\textbf{This is a theory-first modeling project with a proposed pilot validation study.}

\vspace{0.4em}
\textbf{Hypotheses}
\begin{itemize}
\setlength{\itemsep}{0.15em}
\item \textbf{H1:} Moderate segmentation will outperform no segmentation when materials are dense.
\item \textbf{H2:} Excessive segmentation will hurt transfer more than recall.
\item \textbf{H3:} Integrated representations will outperform split representations by lowering extraneous burden.
\end{itemize}

\vspace{0.2em}
\textbf{Validation Path / Study Design}
\begin{itemize}
\setlength{\itemsep}{0.15em}
\item Use structured learning materials decomposable into elements and relations.
\item Compare segmentation (\texttt{none / moderate / excessive}) and representation (\texttt{integrated / split}) conditions.
\item Measure prior knowledge, recall, transfer, and an optional short load rating.
\item Test whether the model's directional predictions match observed learning outcomes.
\end{itemize}

\end{frame}

\end{document}
